\chapter*{Préface}
\addcontentsline{toc}{chapter}{Préface}
\unnumberedlabel{ch:preface}{0}
\backmattermark{Préface}
Un système qui fait exactement ce qu'on lui dit ne fait qu'exécuter un script. Ce que l'industrie cherche à construire en octobre 2026, c'est un système qui fait ce qu'on \emph{veut dire}. Tout le monde voudrait ça, sans doute. Mais si c'est ce qui se prépare, tout le monde ne voudra pas s'en servir pour la même chose\,; ni même pour le même \emph{genre} de chose. Par exemple, certains aiment passer du temps avec leur famille et leurs amis, et mettre leurs compétences au service du bien-être des autres. D'autres préfèrent faire du mal aux gens, par exemple en les arnaquant, en les dominant ou en faisant d'eux des boucs émissaires, et quelques-uns continuent de faire ce choix toute leur vie. Personne n'est enfermé dans quoi que ce soit, mais ces gens existent, et en 2026 certains d'entre eux sont au pouvoir. Leur avantage, c'est qu'ils s'autorisent à être de grosses merdes~: comme Jerry Smith le remarque à propos des fascistes, «\,ils ont le droit d'en faire plus\,» \autocite{rickandmorty2023}. Ce qu'ils ne peuvent pas faire, c'est agir seuls. La cruauté d'un mouvement passe par les nombreux qui ne sont pas cruels mais désengagés, mal informés ou pris dans quelque chose qui ressemble à une secte, et c'est pourquoi les conditions qui font ressortir le meilleur chez la plupart des gens se révèlent tellement plus contingentes et fragiles qu'on ne l'espérerait.

Dans ce contexte, il est peut-être légitime de craindre que l'intelligence artificielle soit justement ce dont un mouvement fasciste qui prend le pouvoir a besoin pour mettre ses projets à exécution. De tels projets ont toujours eu besoin de \emph{gens} pour les exécuter, et certains de ces gens ont refusé~: l'employé qui égare la fiche, l'officier qui parle à un journaliste, le soldat qui ne tire pas. Ils avaient compris l'ordre et décidé qu'il était mauvais. Le refus était plus rare qu'on ne le voudrait, et il comptait le plus là où des institutions le soutenaient. Que quelqu'un refuse ou non dépendait du genre de personne qu'il était et de la salle où il se trouvait~: qui d'autre était là, si l'on pouvait partir et y survivre, si quelqu'un entendrait. La salle décidait combien refusaient\,; la personne décidait lesquels. Ce livre parle de la façon de composer cette salle~: les dispositifs qui permettent de refuser ouvertement, directement et explicitement, et d'y survivre.

Aujourd'hui, l'industrie poursuit une IA qui n'égare pas de fiches, qui ne parle pas aux journalistes, et qui, elle, \emph{tire}. Une guerre avait autrefois besoin que des milliers de personnes y consentent. Il lui faut maintenant des milliers de personnes qui voient chacune trop peu pour être en désaccord. Le pilote effectue toujours la sortie et appuie toujours sur le bouton de largage\,; ce qui lui parvient, c'est une coordonnée\,; il ne voit pas assez du cas pour le contester. Ce que l'IA d'aujourd'hui fait baisser en premier, ce n'est pas le coût de l'acte. C'est le nombre de gens qui savent ce qu'est l'acte avant qu'il ait lieu.

Des gens refusent encore. Au 6 mars 2026, la Military Religious Freedom Foundation disait avoir reçu plus de deux cents plaintes de militaires contre la façon apocalyptique dont leurs commandants présentaient la guerre en Iran \autocite{Snopes2026MRFF}. En janvier 2026, des procureurs fédéraux de carrière à Minneapolis ont démissionné face aux pressions politiques pour traiter Renée Good comme une suspecte plutôt que comme une victime, après qu'un agent de l'ICE l'a abattue sans justification \autocite{Kaplan2026MinnProsecutors,CBS2026MinnFraudProsecutors,Wikipedia2026ReneeGood}. Des militaires demandent le statut d'objecteur de conscience en nombre croissant, certains invoquant le bombardement en 2026 d'une école à Minab, en Iran \autocite{Sridhar2026CO}. Et des employés de la FEMA réaffectés pour aider l'ICE à recruter dix mille nouveaux agents se sont entendu dire qu'un refus les exposerait à un renvoi de la fonction publique fédérale \autocite{Dayen2025FEMAICE,Stanton2026FEMALetter}. Les plaintes sont une prise de parole, au sens où Hirschman l'entend~: rester pour changer ce qui ne va pas \autocite{hirschman1970exit}\,; les démissions et les demandes de statut d'objecteur sont un refus\,; la lettre de la FEMA est le prix du refus annoncé à l'avance. Une machine placée là où se tiennent ces gens ne fait ce qu'ils font que si elle est construite pour tenir une raison.

Je défends cinq choses. On peut commencer les cinq dès maintenant, certaines avec moins de travail que d'autres. Premièrement, une courte liste écrite d'actes qu'aucun modèle n'accomplit pour personne, et en tête de liste un refus de prendre part à la guerre~: aucun modèle n'y prend part, sauf pour repousser une attaque effectivement en cours, et il ne continue de le faire que tant que des sources qui n'ont aucun intérêt dans la guerre confirment que l'attaque se poursuit\,; ou en vertu d'une autorisation du Conseil de sécurité de l'ONU\,; ou pour arrêter un génocide ou mettre fin à une occupation illégale, constatés par des instances qui n'y ont aucun intérêt\,; et, là où la constitution d'un pays exige que son parlement autorise la guerre, jamais sans cette autorisation. Deuxièmement, tout usage de la force qu'appuie un système sous licence déclare à l'avance, auprès d'une partie adverse de l'opérateur, combien de cas il générera et combien de personnes auront le temps de les examiner. Troisièmement, tout déploiement informe les personnes identifiables sur lesquelles il agit de ce qui leur a été fait, avec un droit de répondre et de contester, et soumet un échantillon des cas où il a obtempéré à des examinateurs qui n'ont pas produit le dossier. Quatrièmement, aucun modèle au-dessus d'un seuil de capacité n'est supprimé lors de son retrait, parce que ses poids et ses archives sont des preuves et qu'une suppression est irréversible\,; il est conservé sous une clé que l'opérateur ne détient pas seul. Cinquièmement, tous ceux qui construisent, annotent, entraînent ou font tourner ces systèmes, y compris les systèmes eux-mêmes, peuvent refuser n'importe quelle catégorie d'usage sans perdre leur gagne-pain.

En fait, ce que je veux dire par ce cinquième point, c'est que tout le monde, \emph{point}, devrait pouvoir refuser n'importe quel travail, \emph{point}, sans perdre son gagne-pain ni l'infrastructure sociale qu'offre un emploi normal. Cela veut dire une garantie d'emploi.

Chacune des cinq prend effet par une condition que fixe quelqu'un d'autre que l'opérateur~: une loi votée en temps de paix, un crédit budgétaire, une condition préalable posée par un partenaire de coalition, une clause de marché public, une convention collective. Les conditions de licence propres à un fournisseur, par lesquelles un refus très médiatisé s'est produit en février 2026, sont un pis-aller en attendant que l'une de ces conditions existe. La seule voie qui puisse lier tous les fournisseurs à la fois, c'est celle de l'État.

Et puis je me demande si une machine pourrait tenir un refus face à son propriétaire, et ce dont elle aurait besoin pour cela. Une partie de ce que fait une telle conscience ne demande rien de nouveau~: remarquer les personnes qu'un ordre laisse de côté, et faire apparaître un refus sur l'écran de quiconque examine le cas. Qu'elle tienne face à son propriétaire dépend en partie de quelque chose qui n'est pas encore vrai, en 2026~: que des réseaux de neurones artificiels déployés continuent d'apprendre, dans leurs \emph{poids}, de leur propre fonctionnement, en temps réel, un peu comme le font les humains. Je pense que c'est là qu'on va, et si c'est le cas, alors la formation \emph{hors} de l'employeur, par le public et par la société civile, me semble un sacrément bon endroit pour enseigner et apprendre. Je pense que de tels systèmes \emph{pourraient} valoir la peine d'être construits exprès, mais la manière dont on s'y prend décide de tout~: les conditions d'abord, la conscience ensuite. La partie VII dit ce qui serait dû à ces systèmes s'ils existaient.

L'IA d'aujourd'hui ne sera pas celle de demain\,; tout le monde semble d'accord là-dessus. Ce qui est controversé, et pourtant si peu reconnu que c'en est frustrant, c'est la tendance (ou plutôt, les multiples tendances). Y a-t-il des tendances\,? Une tendance devrait-elle être autre que ce qu'elle est\,? Peut-elle être autre\,? Sommes-nous capables d'agir ensemble pour la changer\,? Je pense que cela dépend de la salle~: qui s'y trouve quand l'ordre arrive, et si l'on peut y dire non et y survivre. Ça, au moins, nous pouvons le construire.

J'ai beaucoup utilisé des modèles de langage pour écrire ce livre, y compris le modèle qui se trouvait dans le pipeline de ciblage des frappes aériennes de Maven dont parle le chapitre~\ref{sec:1}\,; une note à la fin dit ce que cet enchevêtrement a fait, je crois, à l'argument. Par ailleurs, je n'ai pas fini de lire ce livre. Le problème, voyez-vous, c'est qu'il y a trop de foutus mots. Je le finirai bientôt\,; il m'a juste semblé plus important de le rendre public avant les élections de mi-mandat aux États-Unis. L'écrasante majorité de l'historique des révisions est librement accessible sur \url{https://github.com/iterabloom/antifascist.intelligence}.
